\documentclass[../../main.tex]{subfiles}

\begin{document}

\section{Gaussian Elimination}

When you solve a linear system, you probably multiply constants
and add or subtract equations to cancel variables. That is exactly
what we will be doing in Gaussian Elimination, but with matrices.
First, define how we solve a linear system in a more rigorous way.

\begin{definition}{}{}
For a system of linear equations, performing the following operations
will result in a linear system with an identical solution.
\begin{enumerate}[noitemsep]
\item
Change the order of linear equations in which they are listed
\item
Multiply a nonzero constant, or nonzero scalars, to an equation
\item
Add equations from the system to form a new equation
\end{enumerate}
\end{definition}

Now before we connect this idea with matrices, let's define a matrix
that we will use with the property of a linear system above.

\begin{definition}{}{}
An \textbf{augmented matrix} for a linear system $Ax = b$ is a
partitioned matrix in the form $[A \mid b]$.
\end{definition}

Consider the following system for an example.
\begin{align*}
    x + 2y + 3z &= 4 \\
    2x + 3y + 4z &= 5 \\
    3x + 4y + 5z &= 6
\end{align*}
In the linear system above, the augmented matrix will be the following
matrix.
\[
\begin{bNiceArray}{ccc|c}
    1 & 2 & 3 & 4 \\
    2 & 3 & 4 & 5 \\
    3 & 4 & 5 & 6
\end{bNiceArray}
\]
Now that we have an augmented matrix in mind, let's take a look at
how the three properties that we discussed apply to matrices.

\begin{definition}{}{}
For an augmented matrix for a linear system, performing the following
operations will result in a linear system with an identical solution
as the initial system.
\begin{enumerate}[noitemsep]
\item[$R_1$.]
Rearrangement of the rows in the augmented matrix
\item[$R_2$.]
Multiply a nonzero scalar to a row in the augmented matrix
\item[$R_3$.]
Add a row multiplied by a scalar with another to replace a row with
new row
\end{enumerate}
Such operations are referred to as \textbf{elementary row operations}.
\end{definition}

If you compare the operations one by one, you can notice that the
two sets of properties refer to the same property. Now with the
set of properties of augmented matrices in mind, here are the steps
for the Gaussian Elimination algorithm.

\begin{theorem}{}{Gaussian Elimination}
Execute the following steps in order for Gaussian Elimination.
\begin{enumerate}[noitemsep]
\item
Construct the augmented matrix.
\item
Utilize the elementary row operations to transform the augmented
matrix into REF.
\item
Convert the augmented matrix in REF into equations of a transformed
linear system, referred to as the \textbf{derived set}.
\item
Back substitute to solve the system.
\end{enumerate}
\end{theorem}

Let's take a look at a few examples.

\begin{problem}{}{}
Find solution(s) to the following system of linear equations if any.
\begin{align*}
    x + 2y + z &= 1 \\
    3x - y + z &= 6 \\
    x + y - z &= -2
\end{align*}
\end{problem}
\begin{solution}
First, the augmented matrix can be found from the system. Consider
the matrix below.
\[
\begin{bNiceArray}{ccc|c}
    1 & 2  & 1  & 1 \\
    3 & -1 & 1  & 6 \\
    1 & 1  & -1 & -2
\end{bNiceArray}
\]
Using the elementary row operations, the augmented matrix can be
transformed as the following.
\[
\begin{bNiceArray}{ccc|c}
    1 & 2  & 1  & 1 \\
    3 & -1 & 1  & 6 \\
    1 & 1  & -1 & -2
\end{bNiceArray}
\rightarrow
\begin{bNiceArray}{ccc|c}
    1 & 2  & 1  & 1 \\
    0 & -4 & 4  & 12 \\
    1 & 1  & -1 & -2
\end{bNiceArray}
\rightarrow
\begin{bNiceArray}{ccc|c}
    1 & 2  & 1  & 1 \\
    0 & -4 & 4  & 12 \\
    0 & -1 & -2 & -3
\end{bNiceArray}
\]
\[
\rightarrow
\begin{bNiceArray}{ccc|c}
    1 & 2  & 1  & 1 \\
    0 & -1 & 1  & 3 \\
    0 & 1  & 2  & 3
\end{bNiceArray}
\rightarrow
\begin{bNiceArray}{ccc|c}
    1 & 2  & 1  & 1 \\
    0 & -1 & 1  & 3 \\
    0 & 0  & 3  & 6
\end{bNiceArray}
\rightarrow
\begin{bNiceArray}{ccc|c}
    1 & 2  & 1  & 1 \\
    0 & 1  & -1 & -3 \\
    0 & 0  & 1  & 2
\end{bNiceArray}
\]
Because the augmented matrix is now in REF, the transformed
linear system can be rewritten.
\begin{align*}
    x + 2y + z &= 1 \\
    y - z &= -3 \\
    z &= 2
\end{align*}
Back substituting, it is evident that $z = 2$ and $y = -1$.
Moreover, $x = 1 + 2 - 2 = 1$. Therefore, the triple $(1, -1, 2)$
is the solution to the linear system.
\end{solution}

Try substituting to the original system! You can notice that the
solution is valid. Below is the second example of using
Gaussian Elimination.

\begin{problem}{}{}
Find solution(s) to the following system of linear equations if any.
\begin{align*}
    x + 2y + 3z &= 4 \\
    y + 3z &= 5 \\
    x + y &= 1
\end{align*}
\end{problem}
\begin{solution}
First and foremost, the augmented matrix can be found.
\[
\begin{bNiceArray}{ccc|c}
    1 & 2 & 3 & 4 \\
    0 & 1 & 3 & 5 \\
    1 & 1 & 0 & 1
\end{bNiceArray}
\rightarrow
\begin{bNiceArray}{ccc|c}
    1 & 2  & 3  & 4 \\
    0 & 1  & 3  & 5 \\
    0 & -1 & -3 & -3
\end{bNiceArray}
\rightarrow
\begin{bNiceArray}{ccc|c}
    1 & 2 & 3 & 4 \\
    0 & 1 & 3 & 5 \\
    0 & 0 & 0 & 2
\end{bNiceArray}
\]
Writing the transformed augmented matrix into a system of linear
equations, the following system is obtained.
\begin{align*}
    x + 2y + 3z &= 4 \\
    y + 3z &= 5 \\
    0 &= 2
\end{align*}
Notice that $0 = 2$ can never be true. Therefore, no solution exists.
\end{solution}

We could double check our answer without using matrices. Notice that
$3z - x = 4$ and $x = 3z - 4$. Similarly, $y = 5 - 3z$. Substituting
to the first equation, $3z - 4 + 10 - 6z + 3z = 4$, or $6 = 4$,
which is not true. Therefore, our results align! Below is the third
example of using the algorithm.

\begin{problem}{}{}
Find solution(s) to the following system of linear equations if any.
\begin{align*}
    x - y - 2z &= 3 \\
    2x + y - 7z &= 3 \\
    x + y - 4z &= 1
\end{align*}
\end{problem}
\begin{solution}
Consider the following augmented matrix.
\[
\begin{bNiceArray}{ccc|c}
    1 & -1 & -2 & 3 \\
    2 & 1  & -7 & 3 \\
    1 & 1  & -4 & 1
\end{bNiceArray}
\]
Utilizing Gaussian Elimination, the augmented matrix can be
transformed as following.
\begin{align*}
\begin{bNiceArray}{ccc|c}
    1 & -1 & -2 & 3 \\
    2 & 1  & -7 & 3 \\
    1 & 1  & -4 & 1
\end{bNiceArray}
&\rightarrow
\begin{bNiceArray}{ccc|c}
    1 & -1 & -2 & 3 \\
    2 & 1  & -7 & 3 \\
    0 & 2  & -2 & -2
\end{bNiceArray}
\rightarrow
\begin{bNiceArray}{ccc|c}
    1 & -1 & -2 & 3 \\
    0 & 1  & -1 & -1 \\
    2 & 1  & -7 & 3
\end{bNiceArray} \\
&\rightarrow
\begin{bNiceArray}{ccc|c}
    1 & -1 & -2 & 3 \\
    0 & 1  & -1 & -1 \\
    0 & 3  & -3 & -3
\end{bNiceArray}
\rightarrow
\begin{bNiceArray}{ccc|c}
    1 & -1 & -2 & 3 \\
    0 & 1  & -1 & -1 \\
    0 & 0  & 0  & 0
\end{bNiceArray}
\end{align*}
Rewriting the system, the following equations are obtained.
\begin{align*}
    x - y - 2z &= 3 \\
    y - z &= -1
\end{align*}
Notice that $y = z - 1$ and $x = y + 2z + 3 = z - 1 + 2z + 3 =
3z + 2$. Because there exists infinitely many $z$, the system
has infinitely many solutions.
\end{solution}

When writing such solutions, it could be written as the following.
\[
x =
\begin{bmatrix}
    3z + 2 \\
    z - 1 \\
    z
\end{bmatrix}
\]
For the next section, let's discuss the inverse of a matrix and its
application in linear systems.

\end{document}


